\sectionCentralizada{RESUMO}

\vspace{-0.3cm}

\novoBlocoPreTextual{
    \noindent 
    As dúvidas dos estudantes constituem informações relevantes para a prática docente, porém nem sempre chegam ao conhecimento dos professores, seja pela resistência dos alunos em manifestá-las, seja pela utilização de ferramentas alternativas de apoio à aprendizagem. Nessa perspectiva, este trabalho teve como objetivo desenvolver e avaliar um protótipo web de assistente educacional baseado em Modelos de Linguagem de Grande Escala (LLM, do inglês Large Language Models), e Geração Aumentada por Recuperação (RAG, do inglês Retrieval-Augmented Generation), destinado à classificação e organização de dúvidas discentes conforme os tópicos do plano de ensino do professor, disponibilizando essas informações em um painel para acompanhamento dos conteúdos questionados pelos estudantes. A pesquisa caracteriza-se como aplicada, quantitativa e exploratória, combinando pesquisa bibliográfica e experimental. O protótipo foi desenvolvido utilizando \textit{Next.js}, \textit{Supabase}, o modelo de linguagem \textit{Qwen2.5-7B-Instruct}, executado localmente, e o modelo de \textit{embeddings} \textit{multilingual-e5-base}. A plataforma permite que estudantes enviem dúvidas por meio de um \textit{chat} e recebam respostas fundamentadas nos materiais didáticos, enquanto os professores visualizam, em um painel, a distribuição das dúvidas por tópico, acompanhada de resumos individuais e perguntas originais sem identificação dos estudantes. Para a avaliação experimental, foram utilizadas 1.416 perguntas de Estrutura de Dados, provenientes de um conjunto de dados público e traduzidas para o português, distribuídas em 12 tópicos. A classificação realizada pela plataforma foi comparada aos tópicos de referência e obteve acurácia de 83,33\%, precisão macro de 73,85\%, revocação macro de 84,76\% e macro F1-score de 74,90\%. Os seis tópicos com mais dúvidas foram os mesmos na referência e na plataforma. Os resultados indicam que a combinação de LLM e RAG permite organizar automaticamente as dúvidas dos estudantes conforme os tópicos do plano de ensino do professor, dispensando a categorização manual de cada dúvida. Essa organização evidencia os tópicos que concentram mais dúvidas e oferece ao docente uma visão das dificuldades da turma em cada conteúdo, servindo como apoio ao planejamento das aulas.
    \noindent 
}
    
\vspace{0.5cm}
\newcommand{\palavrasChave}{\noindent 
        Palavras-chave: inteligência artificial na educação; LLM; RAG; classificação de dúvidas discentes; prática docente.}
\begin{flushleft}
    \palavrasChave
\end{flushleft}
        


% Obs.: as palavras-chave devem ser escolhidas com bastante rigor, pois devem representar adequadamente os principais temas abordados pela pesquisa
